\chapter{饮用水卫生}
\setcounter{chapter}{6}
\chapterintro{
掌握介水传染病的概念及流行特点；掌握集中式供水的净化与消毒原理（混凝沉淀、过滤、氯化消毒）；掌握生活饮用水水质标准的制定原则；了解饮用水污染的应急处理原则。
}

\section{饮用水污染与疾病}

\subsection{介水传染病}

\begin{defbox}
\textbf{介水传染病（Water-borne Infection）}：由于饮用或接触受病原体污染的水而引起的一类传染病。
\end{defbox}

\begin{keybox}
\textbf{介水传染病的流行特点：}
\begin{enumerate}[leftmargin=*]
    \item 流行呈\imp{暴发型}，多数病人集中在同一潜伏期内
    \item \imp{病例分布与供水范围一致}，绝大多数患者有饮用同一水源的历史
    \item 多发生在\imp{农村地区}，饮生水者多见
    \item 一旦对污染水源采取有效措施后，疾病流行能\imp{迅速得到控制}
    \item 流行病学调查可找到水源污染途径，水质细菌学检查有异常改变
\end{enumerate}
\textbf{常见疾病：}霍乱、伤寒、副伤寒、痢疾、甲型肝炎等。
\end{keybox}

\subsection{化学性污染与健康}

\begin{enumerate}[leftmargin=*]
    \item \imp{氰化物}：剧毒，抑制细胞色素氧化酶，造成细胞内窒息。
    \begin{itemize}
        \item 毒作用机制：①CN$^-$在胃酸作用下→氰氢酸→游离氰离子+细胞色素氧化酶Fe$^{3+}$→氰化高铁细胞色素氧化酶→Fe$^{3+}$失去传递电子能力→中断呼吸链→阻断细胞内氧化代谢→细胞窒息死亡；②氰化物在体内酶作用下生成硫氰酸盐，抑制甲状腺聚碘功能，妨碍甲状腺激素合成，造成甲状腺肿大
        \item 急性中毒分4期：前驱期（刺激症状、呼气带苦杏仁味）→呼吸困难期（胸闷、心悸）→惊厥期（强制性痉挛）→麻痹期（肌肉松弛、反射消失），重者昏迷死亡
        \item 慢性中毒：神经衰弱综合征、神经细胞退行性变
        \item 饮用水中氰化物限值：0.05 mg/L
    \end{itemize}
    \item \imp{硝酸盐}：在体内可还原为亚硝酸盐，导致\keyword{高铁血红蛋白血症}（蓝婴综合征——婴幼儿特别是6个月以内婴儿血中10\%左右血红蛋白转变为高铁血红蛋白，出现发绀等缺氧症状）；与胺类生成\keyword{亚硝胺}（2A类致癌物——消化道癌症）。\textbf{婴儿为敏感人群的原因：}①婴儿胃内pH偏高（5--7）；②婴儿血红蛋白较敏感；③缺乏使高铁血红蛋白还原的酶系统；④单位体重摄水量为成年人3倍
    \item \imp{隐孢子虫}：寄生在人和动物体内，以卵囊形式从粪便排出。卵囊特点：体积小（直径4--6 $\mu$m）、感染剂量低（0.014--0.5个/L）、能抵抗常用饮水消毒剂、普通氯化消毒不能完全去除、但不耐热。\imp{贾第鞭毛虫}也是寄生于人类和动物肠道的原生动物，一般消毒方法很难将其全部杀死。两者均可通过自来水厂过滤去除
    \item \imp{饮水消毒副产物（DBPs）}：氯化消毒过程中氯与水中天然有机物（有机前体物：腐殖酸、富里酸、藻类及其代谢物、蛋白质等）反应生成。包括三卤甲烷（THMs）、卤乙酸（HAAs）等，具有"三致"效应。影响因素：有机前体物含量、加氯量、溴离子浓度、pH（THMs随pH↑而↑，HAAs随pH↑而↓）、温度、接触时间。"Mutagen X"（MX，3-氯-4-二氯甲基-5-羟基-2(5H)呋喃酮）在氯化消毒饮水中浓度很低（数ng/L至数十ng/L），却是\imp{迄今为止最强的诱变物之一}。NDMA（N-二甲基亚硝胺）为2A类致癌物。流行病学研究在膀胱癌和直肠癌显示与DBPs有一致性
\end{enumerate}

\subsection{二次供水污染}

\begin{defbox}
\textbf{二次供水（Secondary Water Supply）}：用水单位将来自市政供水或自备井的生活饮用水贮存于储水池或水箱中，再通过机械加压或凭借高层建筑的自然压差，二次输送至用户的供水系统。我国高层建筑5层以上用水都需要二次供水系统。
\end{defbox}

\begin{tipbox}
\textbf{高层建筑二次供水污染的原因：}
\begin{itemize}[leftmargin=*]
    \item 储水设备问题：储水箱材料不合格（材质老化、侵蚀、涂料脱落）；水箱容积过大，水滞留时间过长，余氯耗尽，微生物繁殖；储水箱设计不合理（出水口高出水箱底平面形成死水）；未及时维修保养、未定期清洗消毒
    \item 管网问题：管道错误连接（上水管设在污水管下面、溢水管与污水管相连无防倒灌措施）；管网破裂渗漏、降压和停水维修时外界污物渗入；管网陈旧老化，水管材质和涂料脱落析出
\end{itemize}
\textbf{健康危害：}以生物性污染为主，通常引起介水肠道传染病；若储水箱材料中有害物质（铅、镉、砷、汞等）含量或溶出过高，可导致慢性危害。
\end{tipbox}

\section{生活饮用水水质标准}

\begin{keybox}
\textbf{制定原则：}
\begin{enumerate}[leftmargin=*]
    \item \imp{流行病学上安全}——不含病原微生物
    \item \imp{化学组成对人体无害}——不引起急慢性中毒及远期危害
    \item \imp{感官性状良好}——色、臭、味、浑浊度满足要求
    \item \imp{经济技术上可行}——标准限值可在现实条件下实现
\end{enumerate}
\end{keybox}

常规检测项目：感官性状和一般化学指标、毒理学指标、细菌学指标、放射性指标。

\begin{tipbox}
\textbf{生活饮用水水质常规指标分类（GB 5749-2022）：}
\begin{itemize}[leftmargin=*]
    \item \imp{微生物学指标}：总大肠菌群（不得检出/100mL）、耐热大肠菌群（不得检出/100mL）、大肠埃希菌（不得检出/100mL）、菌落总数（$\leq$100 CFU/mL）。\textbf{耐热大肠菌群/粪大肠菌群}是判断饮用水是否受粪便污染的重要指标
    \item \imp{消毒剂指标}：游离氯出厂水$\geq$0.3 mg/L（接触30min），末梢水$\geq$0.05 mg/L；末梢水游离氯余量可作为输配水过程中有无再次污染的信号
    \item \imp{毒理学指标}：砷（I类致癌物）、镉、铬（六价）、铅、汞、硒、氰化物、氟化物（$\leq$1 mg/L）、硝酸盐、三氯甲烷、四氯化碳、溴酸盐、甲醛、亚氯酸盐、氯酸盐
    \item \imp{放射性指标}：总$\alpha$放射性、总$\beta$放射性
    \item \imp{感官性状和一般化学指标}：色度、浑浊度（限值1 NTU）、臭和味、肉眼可见物、pH（6.5--8.5）、铝、铁（$\leq$0.3 mg/L）、锰（$\leq$0.1 mg/L）、铜、锌、氯化物、硫酸盐、溶解性总固体、总硬度（$\leq$450 mg/L）、耗氧量、挥发酚类、阴离子合成洗涤剂
\end{itemize}

\textbf{非常规指标：}微生物学2项（贾第鞭毛虫、隐孢子虫）；毒理学39项（农药、除草剂、苯化合物、微囊藻毒素-LR等）；感官性状3项。
\end{tipbox}

\begin{tipbox}
\textbf{部分水质标准制定依据：}
\begin{itemize}[leftmargin=*]
    \item \imp{浑浊度}：限值1度。10度时人即可觉察浑浊；降低浑浊度对去除有害物质、提高消毒效果有积极作用
    \item \imp{pH}：限值6.5--8.5。过低腐蚀管道，过高析出溶解性盐类并降低氯消毒效果
    \item \imp{总硬度}：限值450 mg/L（以CaCO$_3$计）。硬度过高可引起胃肠功能暂时性紊乱
    \item \imp{铁}：限值0.3 mg/L。1 mg/L时有明显金属味，0.5 mg/L时色度可$>$30度
    \item \imp{锰}：限值0.1 mg/L。微量即可使水呈黄褐色；超过0.15 mg/L时使衣物和瓷器出现色斑。锰虽有神经毒性，但因感官影响阈值低于毒理学阈值，按最敏感指标原则归为感官性状指标
    \item \imp{氟化物}：限值1 mg/L。综合考虑1 mg/L时对牙齿的轻度影响和防龋作用，及高氟地区除氟在经济技术上的可行性
\end{itemize}
\end{tipbox}

\section{集中式给水}

\subsection{水源选择原则}

\begin{enumerate}[leftmargin=*]
    \item \imp{水量充足} \item \imp{水质良好} \item \imp{便于防护} \item \imp{技术经济合理}
\end{enumerate}

\subsection{水质净化与消毒}

\begin{keybox}
\textbf{（一）混凝沉淀（Coagulation Precipitation）}

\textbf{原理：}
\begin{enumerate}[leftmargin=*]
    \item \imp{压缩双电层作用}——降低$\zeta$电位
    \item \imp{电中和作用}——中和胶体负电荷
    \item \imp{吸附架桥作用}——高分子混凝剂链状结构吸附多个胶体颗粒
\end{enumerate}

\textbf{混凝剂比较：}

\begin{table}[H]
\centering
\caption{常用混凝剂比较}
\begin{tabularx}{\textwidth}{lXX}
\hline
\textbf{特性} & \textbf{铝盐（明矾、硫酸铝）} & \textbf{铁盐（三氯化铁、硫酸亚铁）} \\
\hline
适宜pH & 窄（6.5--7.5） & 宽（5--9） \\
絮凝体 & 慢且松散（低温时更明显） & 快、大而紧密 \\
腐蚀性 & 小 & 大，易潮湿 \\
对水质影响 & 无不良影响 & 处理后含铁量高 \\
低温低浊水效果 & 较差 & 较好 \\
\hline
\end{tabularx}
\end{table}

\imp{聚合氯化铝（PAC）}：适用pH范围宽、腐蚀性小、成本低、絮体形成快、沉淀效果好、用量少。

\imp{聚丙烯酰胺（PAM）}：高分子混凝剂，常作为助凝剂使用。

\textbf{助凝剂}：①调节混凝条件（原水碱度不够加石灰；用氯将亚铁氧化为高铁）；②改善絮凝体结构（铝盐产生的絮凝体细小松散时，用聚丙烯酰胺或活性硅酸等助凝）。

\textbf{影响因素：}水温；水的pH（对金属盐类混凝剂影响较大，对高分子混凝剂影响较小）；水中微粒的性质和含量；水中有机物和溶解盐含量；混凝剂种类和用量；投加方法、搅拌强度和反应时间等。

\textbf{（二）过滤（Filtration）}

\textbf{原理：}浑水通过滤料层，悬浮颗粒和微生物被截留（筛除作用+接触凝聚作用）。工作周期：成熟期→过滤期→清洗期。

\textbf{滤料要求：}本身无毒、化学性质稳定；不能被微生物利用和分解；有良好机械强度；颗粒粒度均匀，有适当孔隙率。常用滤料：石英砂、无烟煤、木炭、活性炭、磁铁砂。

\textbf{影响因素：}滤层厚度和粒径、滤速和滤池类型、进水水质（浑浊度、色度、有机物、藻类）。

\textbf{（三）消毒（Disinfection）}

\textbf{理想消毒剂的选择标准：}①能迅速灭活病原微生物（细菌+病毒+真菌+寄生虫）；②在使用的浓度水平无毒；③便于控制和检测；④有剩余消毒剂；⑤对水感官性状无不良影响；⑥副产物对健康影响小且可预防或消除；⑦经济技术上可行。

\textbf{氯化消毒（Chlorination Disinfection）}

\textbf{原理：}含氯制剂加入水中生成\imp{次氯酸（HOCl）}，HOCl为中性小分子，强渗透性，强氧化剂，损害细胞膜使通透性增加，蛋白质、RNA、DNA等释放，影响多种酶系统（细菌：干扰糖代谢；病毒：对核酸致死性损害）。

\begin{defbox}
\textbf{有效氯}：含氯化合物分子团中氯的价数 $>$ -1，是含氯化合物中有杀菌能力的有效成分。

\textbf{需氯量}：因杀菌、氧化水中有机物和还原性无机物需要的总氯量。

\textbf{余氯（Residual Chlorine）}：加氯量 $>$ 需氯量，使在氧化和杀菌后还能剩余的有效氯量。是评价氯化消毒效果是否达标的简便指标。
\begin{itemize}
    \item \imp{游离性余氯}：HOCl + OCl$^-$，消毒效果好。接触30 min，含量0.3--0.5 mg/L
    \item \imp{化合性余氯}：NH$_2$Cl + NHCl$_2$，消毒效果略差。接触1--2 h，含量1--2 mg/L
\end{itemize}
\end{defbox}

\textbf{HOCl vs OCl$^-$}：HOCl的杀菌效率比OCl$^-$高80倍。pH $<$ 6时HOCl $>$ 95\%；pH = 7.5时HOCl = OCl$^-$；pH $>$ 9时OCl$^-$ $\approx$ 100\%。

\textbf{影响因素：}
\begin{enumerate}[leftmargin=*]
    \item \imp{加氯量}——越多效果越好，但易产生DBPs
    \item \imp{接触时间}——越长效果越好。水温每提高10°C，杀菌率提高2--3倍
    \item \imp{pH值}——pH越低HOCl比例越高，效果越好
    \item \imp{水温}——越高杀菌效果越好 → 冬季需延长接触时间
    \item \imp{浑浊度}——高时细菌附着在悬浮颗粒上，氯不易直接作用于细菌
    \item \imp{微生物种类和数量}——抵抗力：大肠杆菌 $<$ 病毒 $<$ 原虫包囊
\end{enumerate}

\textbf{氯化消毒方法：}
\begin{itemize}[leftmargin=*]
    \item \imp{普通氯化消毒法}：水浊度低、有机物污染轻、基本无氨（$<$0.3 mg/L）无酚时使用。接触时间短，效果可靠，但产生氯化副产物
    \item \imp{氯胺消毒法}：水中加氯后再加氨，生成一氯胺和二氯胺。副产物少，稳定性高，适宜长距离输水；但消毒作用不如次氯酸强，接触时间长
    \item \imp{折点氯消毒法}：加氯量超过折点，形成适量游离氯。消毒效果可靠，能降低臭味、色度及有机物含量；但耗氯多，副产物多
    \item \imp{过量氯消毒法}：水源污染严重或需短时间消毒时，加过量氯使余氯达1--5 mg/L，消毒后用亚硫酸钠/活性炭脱氯
\end{itemize}
\end{keybox}

\subsection{消毒方法比较}

\begin{table}[htbp]
\centering
\caption{常用饮水消毒方法比较}
\begin{tabularx}{\textwidth}{lXXXX}
\hline
\textbf{特性} & \textbf{氯化消毒} & \textbf{二氧化氯} & \textbf{臭氧消毒} & \textbf{紫外线消毒} \\
\hline
消毒效果 & 良好 & 优异 & 优异 & 良好 \\
持续消毒能力 & \imp{有} & 有 & 无 & 无 \\
消毒副产物 & THMs等 & 亚氯酸盐 & 溴酸盐等 & 无 \\
成本 & 低 & 中等 & 高 & 中等 \\
\hline
\end{tabularx}
\end{table}

\subsection{氯化消毒副产物的预防}

\begin{itemize}[leftmargin=*]
    \item 减少水源中有机物前体物质
    \item 优化消毒工艺（控制加氯量、改变加氯点）
    \item 采用替代消毒剂（氯胺、二氧化氯、臭氧）
    \item 深度处理去除已生成的副产物
\end{itemize}

\subsection{饮用水的深度净化}

在市政供水常规处理基础上，再次对水质净化处理以获得优质饮用水。
\begin{itemize}[leftmargin=*]
    \item \imp{物理吸附分离}：活性炭吸附法；膜分离法（微滤MF、超滤UF、纳滤NF、反渗透RO）
    \item \imp{化学氧化技术}：预氧化技术、TiO$_2$光催化技术
    \item \imp{生物预处理技术}
\end{itemize}

\section{分散式给水}

\subsection{包装饮用水}
\begin{itemize}[leftmargin=*]
    \item \imp{纯水}：以市政自来水为原水，经反渗透、电渗析、蒸馏等工艺使水中溶解矿物质及其他有害物质全部去除，电导率$<$10 $\mu$s/cm，浊度$<$1度。\textbf{不建议作生活饮用水长期饮用}
    \item \imp{净水}：以市政自来水为原水，通过吸附、超滤、纳滤去除有害物质而保留原水中溶解性矿物质。\textbf{可作生活饮用水长期饮用}
    \item \imp{天然矿泉水}：储存于地下深处自然涌出或人工采集的未受污染且含偏硅酸、锶、锌等一种或多种微量元素达到限量值的泉水
\end{itemize}

\subsection{直饮水（Direct Drinking Water）}
属分质供水范畴，对自来水深度处理后，通过食品卫生级输水管道送入用户，可直接饮用。

\subsection{涉水产品}
在饮用水生产和供水过程中与饮用水接触的连接止水材料、塑料及有机合成管材、管件、防护涂料、水处理剂、水质处理器及其他新材料和化学物质。

\subsection{传统分散式给水}

\begin{itemize}[leftmargin=*]
    \item \textbf{地下水}（井水）：水质较好，需注意卫生防护
    \item \textbf{地表水}：需经处理后方可饮用
    \item \textbf{雨雪水}：水量无保证，易受大气污染影响
\end{itemize}

\section{饮用水污染的应急处理}

\textbf{水污染事故应急处理原则：}
\begin{enumerate}[leftmargin=*]
    \item 立即停止使用受污染水源
    \item 迅速查明污染源和污染物
    \item 采取紧急净化和消毒措施
    \item 启用备用水源或应急供水
    \item 开展人群健康监测和医学救援
    \item 及时向公众通报信息
\end{enumerate}

\section{复习题}

\begin{enumerate}[leftmargin=*, itemsep=8pt]
    \item \textbf{【名词解释】}介水传染病；饮水消毒副产物（DBPs）
    \item \textbf{【名词解释】}混凝沉淀；氯化消毒
    \item \textbf{【简答题】}简述介水传染病的流行特点。
    \item \textbf{【简答题】}简述生活饮用水水质标准的制定原则。
    \item \textbf{【简答题】}简述集中式供水水源选择的原则。
    \item \textbf{【简答题】}简述混凝沉淀的原理及影响因素。
    \item \textbf{【简答题】}简述影响氯化消毒效果的因素。
    \item \textbf{【简答题】}比较四种常用消毒方法的优缺点。
    \item \textbf{【简答题】}简述如何预防和控制氯化消毒副产物。
    \item \textbf{【简答题】}简述高层建筑二次供水污染的原因。
    \item \textbf{【简答题】}饮用水应急事件的调查内容和处理原则是什么？
    \item \textbf{【非标准答案题】}理想的饮用水应该是什么样的？为什么长期饮用纯净水对健康可能不利？
\end{enumerate}

\subsection*{参考答案要点}

\begin{enumerate}[leftmargin=*, itemsep=5pt]
    \item \textbf{介水传染病}：通过饮用或接触受病原体污染的水而引起的传染病。\textbf{DBPs}：消毒剂与原水中有机物反应生成的具有"三致"效应的副产物。

    \item \textbf{混凝沉淀}：加入混凝剂使水中悬浮物和胶体聚合沉淀的净水方法。\textbf{氯化消毒}：使用含氯制剂生成HOCl杀菌的消毒方法。

    \item 五大特点：暴发型、病例分布与供水范围一致、多见于农村、控制水源后迅速控制、水质细菌学异常。

    \item 四项原则：流行病学安全、化学物质无害、感官性状良好、经济技术可行。

    \item 四项原则：水量充足、水质良好、便于防护、技术经济合理。

    \item 原理：压缩双电层+电中和+吸附架桥。影响因素：pH、水温、混凝剂种类用量、搅拌。

    \item 六大因素：加氯量、接触时间、pH值、水温、浑浊度、微生物种类和数量。

    \item 氯化持续效果好成本低但有副产物；臭氧效果好但无持续能力成本高；紫外线无副产物但无持续能力；二氧化氯效果好有持续能力副产物较少。

    \item 减少前体物质、优化消毒工艺、采用替代消毒剂、深度处理去除已生成DBPs。

    \item 储水池设计不合理、储水时间过长、管理不善、管网老化渗漏。

    \item 调查：污染源、污染物、范围、暴露人群、健康影响。处理：停用、查明、紧急处理、备用水源、健康监测、信息公开。

    \item 理想饮用水应安全卫生、含适量矿物质、感官良好。纯净水缺乏Ca、Mg等矿物质，pH偏酸，长期饮用影响矿物质平衡。
\end{enumerate}

% ----- PPT参考题 -----
\subsection*{参考题（PPT课件补充）}

\begin{enumerate}[leftmargin=*, itemsep=8pt]
    \item \textbf{【单选题】}下列哪种疾病属于介水传染病？（~~~~）
    \begin{multicols}{5}
    A. 台湾油症 \\
    B. 霍乱 \\
    C. 水俣病 \\
    D. SARS \\
    E. 氟骨症
    \end{multicols}
    \item \textbf{【单选题】}针对介水传染病的暴发流行，最有效的防控措施是（~~~~）
    \begin{multicols}{5}
    A. 勤洗手 \\
    B. 戴防护口罩 \\
    C. 隔离感染病人 \\
    D. 饮水净化消毒 \\
    E. 接种疫苗
    \end{multicols}
    \item \textbf{【单选题】}有村民反映当地饮用水烧开后结垢多，CDC取水样检测发现水中溶解性总固体为1872 mg/L，总硬度550 mg/L，大大超过生活饮用水卫生标准。这样的水最可能来自（~~~~）
    \begin{multicols}{5}
    A. 江水 \\
    B. 河水 \\
    C. 湖泊水 \\
    D. 井水 \\
    E. 水库水
    \end{multicols}
    \item \textbf{【单选题】}锰具有神经毒性，微量锰可使水呈黄褐色。在生活饮用水卫生标准中，锰被归类为感官性状和一般化学指标，而不是毒理学指标，这是遵循了环境卫生标准制订中的（~~~~）
    \begin{multicols}{3}
    A. 保障居民不发生急性及慢性危害原则 \\
    B. 最敏感指标的原则 \\
    C. 对人体健康无间接影响原则 \\
    D. 最经济指标的原则 \\
    E. 最切实可行的原则
    \end{multicols}
    \item \textbf{【单选题】}卫生监督执法部门在对某自来水厂进行常规监督时发现，该水厂出厂水游离氯余量为0.5 mg/L，而末梢水中几乎不能检测到游离余氯，提示（~~~~）
    \begin{multicols}{2}
    A. 水厂液氯添加量不足 \\
    B. 自来水厂操作不规范 \\
    C. 管网内出现二次污染 \\
    D. 游离余氯转变为化合性余氯
    \end{multicols}
\end{enumerate}

\subsubsection*{参考答案}
\begin{enumerate}[leftmargin=*, itemsep=5pt]
    \item \textbf{B（霍乱）}。介水传染病指通过饮用或接触受病原体污染的水而引起的传染病，霍乱是典型的介水传染病。台湾油症由多氯联苯引起（食物污染），水俣病为慢性甲基汞中毒（化学性污染），SARS为呼吸道传染病，氟骨症为地球化学性疾病。
    \item \textbf{D（饮水净化消毒）}。介水传染病的传播途径是受污染的水源，最有效的控制措施是切断传播途径——对饮用水进行净化消毒。勤洗手属个人卫生，口罩防呼吸道传播，隔离病人针对传染源，疫苗接种针对易感人群，均不如饮水消毒直接有效。
    \item \textbf{D（井水）}。地下水（井水）溶解土壤矿物质较多，硬度较高，溶解性总固体含量高于地表水。江、河、湖、水库水均属地表水，硬度通常较低。
    \item \textbf{B（最敏感指标的原则）}。锰在低于产生神经毒性的浓度时即可使水着色（感官不良），其感官影响的阈值低于毒理学阈值，按最敏感指标的原则归类为感官性状和一般化学指标。
    \item \textbf{C（管网内出现二次污染）}。出厂水有余氯（0.5 mg/L）说明水厂加氯正常；末梢水无游离余氯说明余氯在管网输送途中被消耗，最常见的原因是管网老化渗漏导致二次污染，细菌等污染物消耗了余氯。
\end{enumerate}

\newpage

% =====================================================================
%  CHAPTER 7: 土壤卫生 【重点章节】
% =====================================================================
\newpage